\itemhead{p.97 练习 · 第2题（教材p.97）}{XJ-B1-C04-PR-P097-Y1020}
\label{XJ-B1-C04-PR-P097-Y1020}
{\footnotesize 湘教版教材p.97 / PDF第104页 / p.97 练习 / 第2题}
\field{原题}{用分数指数幂表示下列各式（式中字母都表示正数）：
\subq{1}$\sqrt[4]{a^3}$；
\subq{2}$\sqrt[4]{(m+n)^3}$；
\subq{3}$\sqrt[5]{(a+2b)^2}$；
\subq{4}$\sqrt[4]{(m-n)^2}\ (m<n)$；
\subq{5}$\sqrt{a^4b^3}$；
\subq{6}$\dfrac{m^2}{\sqrt[3]m}$。}
\field{参考答案}{依次为 $a^{3/4},(m+n)^{3/4},(a+2b)^{2/5},(n-m)^{1/2},a^2b^{3/2},m^{5/3}$；第（4）利用 $m<n$，不能直接写负底数的 $1/2$ 次幂。}
{\small \field{关联知识}{P02 有理数、实数指数幂与运算（教材帮p.144、145 / PDF153、154）}
\field{方法与易错点}{统一为指数幂，分开处理系数与各底数的指数；最终复查定义条件及数值精度。 易错：负指数误作负数；幂的幂指数相加；约分丢禁点；过早四舍五入。}
\field{答案依据}{答案册相关题组：PDF第30页（册内页29）；按当前题库原题逐项重算，不按旧版题号机械对套。}}
